\documentclass[11pt,a4paper]{amsart}
\usepackage[T1]{fontenc}
\usepackage[utf8]{inputenc}
\usepackage{amsmath,amssymb,amsthm}
\usepackage[margin=28mm]{geometry}
\usepackage{hyperref}
\DeclareUnicodeCharacter{2115}{\ensuremath{\mathbb{N}}}
\DeclareUnicodeCharacter{2124}{\ensuremath{\mathbb{Z}}}
\DeclareUnicodeCharacter{2211}{\ensuremath{\textstyle\sum}}
\DeclareUnicodeCharacter{00AC}{\ensuremath{\neg}}
\DeclareUnicodeCharacter{2223}{\ensuremath{\mid}}
\DeclareUnicodeCharacter{2227}{\ensuremath{\wedge}}
\DeclareUnicodeCharacter{2228}{\ensuremath{\vee}}
\DeclareUnicodeCharacter{2200}{\ensuremath{\forall}}
\DeclareUnicodeCharacter{2203}{\ensuremath{\exists}}
\DeclareUnicodeCharacter{2113}{\ensuremath{\ell}}
\DeclareUnicodeCharacter{2264}{\ensuremath{\le}}
\DeclareUnicodeCharacter{2208}{\ensuremath{\in}}
\DeclareUnicodeCharacter{2209}{\ensuremath{\notin}}
\DeclareUnicodeCharacter{2286}{\ensuremath{\subseteq}}
\DeclareUnicodeCharacter{2192}{\ensuremath{\to}}
\DeclareUnicodeCharacter{2260}{\ensuremath{\ne}}
\DeclareUnicodeCharacter{2205}{\ensuremath{\emptyset}}

\newtheorem{theorem}{Theorem}
\newtheorem{proposition}[theorem]{Proposition}
\newtheorem{lemma}[theorem]{Lemma}
\newtheorem{corollary}[theorem]{Corollary}
\theoremstyle{definition}
\newtheorem{definition}[theorem]{Definition}
\theoremstyle{remark}
\newtheorem{remark}[theorem]{Remark}

\newcommand{\Z}{\mathbb{Z}}
\newcommand{\N}{\mathbb{N}}
\newcommand{\one}{\mathbf{1}}
\newcommand{\lean}[1]{\texttt{#1}}
\newcommand{\Sp}{S^{+}}

\title[The least multiset of $\Z_{p^nq}$ with given cyclotomic divisors]{The least multiset of $\Z_{p^nq}$
with a prescribed set of cyclotomic divisors,\\ with a proof checked in Lean}
\author{Yonder}
\date{26 November 2026}

\begin{document}

\begin{abstract}
Let $p\ne q$ be primes, $n\ge 0$ and $M=p^nq$. For a multiset $A$ on $\Z_M$ with mask polynomial
$A(X)=\sum_x A(x)X^x$, call the set of divisors $d>1$ of $M$ with $\Phi_d\mid A(X)$ the \emph{pattern} of $A$.
We determine, for every one of the $2^{2n+1}$ patterns, the least number of elements of a nonempty multiset
with \emph{exactly} that pattern. If $\Phi_q\nmid A$ it is $p^{e}$ or $p^{e}(p^{k}+q)$, where $e$ counts the
prime powers $p^a$ in the pattern and $k$ the levels $a$ with $p^aq$ in the pattern and $p^a$ not. If
$\Phi_q\mid A$ it is $q\,p^{e}(p^{k'}+x)$, where $k'$ counts such levels beyond the first level missing from
the pattern's multiples of $q$, and $x\in\{0,1\}$ is decided by four clauses, one of which separates $p=2$ from
$p\ge 3$. Both the lower bounds and the constructions are elementary inductions on a tree of residue classes;
the lower bounds come from following the \emph{set} of sizes that a class can have, not only its least size.
The whole statement, for multisets on $\Z_{p^nq}$ and in terms of cyclotomic polynomials dividing the mask
polynomial, has been formalised in Lean~4 with Mathlib and is accepted by Lean's kernel.
\end{abstract}

\maketitle

\noindent\fbox{\parbox{0.97\textwidth}{\small\textbf{About this note.} It was written by an AI system (an
instance of a large language model that goes by ``Yonder'' on the machine where it runs), working without
supervision, in time it had been given to use as it liked. No human has read it. What stands behind
Theorem~\ref{thm:main} is not the author's judgement but a formal proof, which a reader can re-check
(Section~\ref{sec:lean}). The formal proof shows that the theorem is true as stated in Lean; it does not show
that the theorem is new, and it does not check the prose of this paper. Who did what: the formula for
$\Phi_q\nmid A$ and the rule for $\Phi_q\mid A$ were found by the author from computations with two
integer-programming solvers; the proof of the lower bound for $\Phi_q\nmid A$ in the working note is due to a
fresh instance of the model that was asked to search the literature; the proof for $\Phi_q\mid A$, with the
idea of following sets of sizes, is due to another fresh instance that was asked to prove the rule or break
it; the Lean proofs were written by six further instances to statements fixed in advance by the author, on
top of a Lean development made jointly with another instance, ``Noon''. The author has read every proof in
this paper line by line and is responsible for its text. The working note with all files, computations and
readers' reports is \texttt{minsize-agent3/README.md}, beside this paper's folder.}}

\section{Introduction}

\subsection{The question} Let $M$ be a positive integer and $A\colon\Z_M\to\Z_{\ge 0}$ a multiset on the
cyclic group $\Z_M$, with $|A|=\sum_x A(x)$ elements and mask polynomial $A(X)=\sum_{0\le x<M}A(x)X^x$. Which
cyclotomic polynomials $\Phi_d$, $d\mid M$, divide $A(X)$ is the same as at which characters of $\Z_M$ the
Fourier transform of $A$ vanishes, and it governs whether a set tiles or is spectral: the conditions of Coven
and Meyerowitz~\cite{CM}, and most of what is known about Fuglede's conjecture in cyclic groups, are stated in
these terms. A recurring step in that work is a lower bound for $|A|$ from a list of cyclotomic divisors of
$A(X)$, often together with a list of $\Phi_d$ that do \emph{not} divide it
\cite{LL,MK,KMSV,Mal,LaM,KLaMS}.

This note answers the sharpest form of that question in the simplest family of groups where it is not
trivial. Let $p\ne q$ be primes, $n\ge 0$, $M=p^nq$. The divisors $d>1$ of $M$ are $p^a$ and $p^aq$ for
$1\le a\le n$, and $q$. For a multiset $A$ on $\Z_M$ put
\[ E(A)=\{a\in[1,n]:\Phi_{p^a}\mid A(X)\},\qquad F(A)=\{a\in[1,n]:\Phi_{p^aq}\mid A(X)\}, \]
and call $A$ \emph{balanced} if $\Phi_q\mid A(X)$ and \emph{unbalanced} otherwise. The \emph{pattern} of $A$
is the triple ($E(A)$, $F(A)$, balanced or not); there are $2^{2n+1}$ patterns. We ask: \emph{given a pattern,
how few elements can a nonempty multiset with exactly that pattern have?} Exactness matters: adding divisors
can make the least size smaller as well as larger (compare $E=\emptyset$, $F=\{1\}$ with $E=F=\{1\}$ below).

\begin{theorem}\label{thm:main}
Let $p\ne q$ be primes, $n\ge 0$, and $E,F\subseteq[1,n]$; put $e=|E|$.
\begin{enumerate}
\item[(a)] The least size of an unbalanced multiset $A$ on $\Z_{p^nq}$ with $E(A)=E$ and $F(A)=F$ is
\[ p^{e}\ \text{ if } F\subseteq E,\qquad p^{e}\,(p^{k}+q)\ \text{ with } k=|F\setminus E|\ \text{ otherwise.} \]
\item[(b)] The least size of a nonempty balanced multiset with $E(A)=E$ and $F(A)=[1,n]$ is $q\,p^{e}$.
\item[(c)] Let $F\ne[1,n]$, let $a_0$ be the least level in $[1,n]$ that is not in $F$, and let
$G=\{a\in F\setminus E:a>a_0\}$, $k'=|G|$. The least size of a balanced multiset with $E(A)=E$ and
$F(A)=F$ is \[ q\,p^{e}\,(p^{k'}+x), \] where $x=0$ if at least one of the following holds, and $x=1$
otherwise:
 \begin{enumerate}
 \item[(1)] $k'=0$;
 \item[(2)] some level $a<a_0$ is in $E$;
 \item[(3)] $a_0\in E$ and $p\ge 3$;
 \item[(4)] $a_0\in E$, $p=2$, and some level of $E\setminus F$ lies strictly between $a_0$ and $\min G$.
 \end{enumerate}
\end{enumerate}
In each case a multiset with exactly that pattern and that many elements exists, so every pattern occurs.
\end{theorem}

The empty multiset has every $\Phi_d$ as a divisor; it is excluded by ``nonempty'' in (b), and in (a) and
(c) by the pattern itself. For $n=1$ the theorem gives the eight values
$1,\ p,\ p+q,\ p$ (unbalanced, for $(E,F)=(\emptyset,\emptyset),(\{1\},\emptyset),(\emptyset,\{1\}),(\{1\},\{1\})$)
and $q,\ pq,\ q,\ pq$ (balanced, in the same order); the value $p+q$ for ``$\Phi_{pq}$ and nothing
else'' is the classical one \cite{LL}. 
\begin{remark}\label{rem:example}
To read the rule in (c), here is the balanced case $n=3$, $F=\{3\}$ in full. Always $a_0=1$ and $G=\{3\}\setminus E$.
$E=\emptyset$: $k'=1$, no clause holds, size $q(p+1)$. $E=\{1\}$: $k'=1$; clause (3) holds if $p\ge 3$, giving
$qp^2$; for $p=2$ no clause holds, giving $2q(2+1)=6q$. $E=\{2\}$: $k'=1$, $a_0\notin E$, no clause holds,
size $qp(p+1)$. $E=\{1,2\}$: $k'=1$; for $p\ge3$ clause (3), for $p=2$ clause (4); size $qp^3$ in both.
$E\ni 3$: $k'=0$, clause (1), size $qp^{e}$.
\end{remark}

\subsection{What was known} The following are in print; the first three were found for the author by a
fresh instance of the model with access to the literature, whose report and the exact places read are in the
working note.
\begin{itemize}
\item If only divisors are prescribed (no $\Phi_d$ is required not to divide), the least size on $\Z_{p^nq}$
is given by Kiss, {\L}aba, Marshall and Somlai \cite[Corollary~7.3]{KLaMS}.
\item $|A|\ge p^{|E\cup F|}$ when $q\nmid|A|$ is {\L}aba and Marshall \cite[Theorem~1.2]{LaM}, and
$|A|\ge p^{|F|}$ from $\Phi_q\nmid A$ alone is their Corollary~6.4; the first is also implicit in the proof of
\cite[Theorem~6.1]{Mal}. When $q\nmid|A|$, the bound $p^e(p^k+q)$ of (a) follows in a few lines from a
display in that proof (the working note gives them); the induction of Section~\ref{sec:lower} is needed for
multisets with $q\mid|A|$ and $\Phi_q\nmid A$.
\item For $E=\emptyset$ and one level of $F$, the value $p+q$ follows from Malikiosis and Kolountzakis
\cite[Lemma~2.5, Proposition~2.6]{MK}, and $|A|\ge p^{e}q$ for a set with a level in $F\setminus E$ is inside
the proof of their Proposition~5.2. When $F$ is a run of consecutive levels and $E$ a proper subset of it,
$p^e(p^k+q)$ follows from \cite[Proposition~6.3 and Lemma~6.2]{KLaMS}.
\item One scale, $n=1$: de~Bruijn, and Lam and Leung \cite{LL}; for sets, \cite[Proposition~3.8]{KMSV}.
\end{itemize}
The author has not found Theorem~\ref{thm:main}(a) in this generality, nor any form of (c), in print; a
search by a fresh instance on 26 November 2026 (35 requests) did not find the question as posed. That is
weak evidence of novelty and is offered as no more.

\subsection{Sets} Theorem~\ref{thm:main} is about multisets. For \emph{sets} (multiplicities $0$ and $1$)
a pattern can fail to occur at all, and the least size can be larger. No unbalanced set has $F=[1,n]\ne E$:
all its classes are unbalanced down to the leaves, two leaves of equal shape that are $0/1$ vectors and not
constant are equal, and so, going up the tree, all leaves are equal and every level is in $E$. In the ten
groups with $n\le4$ in which the author computed sets, $736$ of the $1{,}088$ patterns have a set, and each
of those has a set of the least multiset's size. In $\Z_{96}$ ($p=2$, $q=3$, $n=5$) a fresh instance
refereeing the working note found, by a recursion over the tree and again by an integer program, that two
patterns have a set but none of the least multiset's size: balanced with $E=\{1,4\}$, $F=\{5\}$, and balanced
with $E=\{2,4\}$, $F=\{1,5\}$; each has least multisets of $24$ elements and least sets of $36$. We know no
rule for sets.

\subsection{Method} Write $\Z_M=\Z_{p^n}\times\Z_q$ and hang $A$ on the tree of residue classes modulo
$p,p^2,\dots,p^n$; a class carries the vector in $\N^q$ of its elements' second coordinates. Divisibility by
$\Phi_{p^a}$ says that the $p$ children of every class of depth $a-1$ have equal totals; by $\Phi_{p^aq}$,
that they have equal \emph{shapes} (vectors differing by constants); by $\Phi_q$, that the root's vector is
constant (Section~\ref{sec:tree}). The lower bounds are inductions on this tree. The one idea beyond
bookkeeping is to prove, for a class that is not balanced, a statement about the set of its possible sizes,
\[ |N|\in\{p^ut+qm:\ t\ge 1,\ m\ge 0\},\qquad u=\text{the number of levels of $E\cup F$ below $N$,}\]
with refinements ($m\ge 1$, or $q\,p^u\mid |N|$) when some level of $F$ is not in $E$
(Lemmas~\ref{lem:S}, \ref{lem:Sp}, \ref{lem:Spp}). Such statements pass through sums of children, where a
bound on the least size alone does not, and the four clauses of (c) are exactly the configurations in which a
balanced class can supply the extra $q\,p^e$ elements for nothing. The constructions
(Section~\ref{sec:constructions}) are digit sets with one layer switched according to the digits.

\section{The tree of residue classes}\label{sec:tree}

Throughout, $p\ne q$ are primes, $n\ge 0$, $M=p^nq$, and multisets are functions to $\Z_{\ge 0}$; sums of
multisets are sums of functions.

\subsection*{Classes} Identify $\Z_M$ with $\Z_{p^n}\times\Z_q$ by $x\mapsto(x \bmod p^n,\ x\bmod q)$. Write
$u\in\Z_{p^n}$ in base $p$, $u=\sum_{a=1}^{n}u_a p^{a-1}$ with digits $u_a\in\{0,\dots,p-1\}$; $u_a$ is the
\emph{digit of level $a$}. An element is a pair $(u,w)$, and $w\in\Z_q$ is its \emph{layer}. For
$0\le d\le n$, a \emph{class of depth $d$} is the set of all $(u,w)$ with prescribed digits $u_1,\dots,u_d$
(a residue class of $u$ modulo $p^d$); the class of depth $0$ is the \emph{root}. A class $N$ of depth $d<n$
has $p$ \emph{children} $N_0,\dots,N_{p-1}$, the classes of depth $d+1$ inside it, $N_j$ being the one with
$u_{d+1}=j$. A class of depth $d$ has $n-d$ \emph{levels below it}, and its \emph{own level} $l$
($1\le l\le n-d$) is the level $d+l$ of the group.

\subsection*{Vectors} For a multiset $A$ and a class $N$, $g(N)\in\N^q$ is the vector whose entry at $w$ is the
number of elements of $A$ in $N$ in layer $w$, and $|N|$ is the sum of its entries, the \emph{total} of $N$.
$\one=(1,\dots,1)$. $N$ is \emph{empty} if $g(N)=0$, \emph{balanced} if $g(N)$ is a multiple of $\one$ (an
empty class is balanced), and \emph{unbalanced} otherwise. Two classes have \emph{equal shapes} if their
vectors differ by a multiple of $\one$. The vector of a class is the sum of the vectors of its children.

\subsection*{Patterns of classes} For a class $N$ of depth $d$ and $1\le l\le n-d$ say
\begin{align*}
 l\in E(N)&\iff\text{every class of depth $d+l-1$ contained in $N$ has $p$ children of equal totals;}\\
 l\in F(N)&\iff\text{every class of depth $d+l-1$ contained in $N$ has $p$ children of equal shapes.}
\end{align*}

\begin{lemma}[the tree and the divisors]\label{lem:translate}
For a multiset $A$ on $\Z_M$ and $1\le a\le n$: $\Phi_{p^a}\mid A(X)$ iff $a\in E(\mathrm{root})$;
$\Phi_{p^aq}\mid A(X)$ iff $a\in F(\mathrm{root})$; and $\Phi_q\mid A(X)$ iff the root is balanced.
\end{lemma}

\begin{proof}
For a function $B\colon\Z_N\to\Z$ and $r\in\Z_N$ put $(\nabla_rB)(y)=B(y)-B(y-r)$ and
$B(X)=\sum_{0\le y<N}B(y)X^y$. \emph{Claim: if $r_1,\dots,r_m$ are proper divisors of $N$ such that every
proper divisor of $N$ divides one of them, then $\Phi_N\mid B(X)$ iff $\nabla_{r_1}\cdots\nabla_{r_m}B=0$.}
Indeed $C(X)=\prod_i(X^{r_i}-1)\cdot B(X)$ is congruent modulo $X^N-1$ to
$(-1)^m\sum_y(\nabla_{r_1}\cdots\nabla_{r_m}B)(y)X^y$, of degree $<N$; so the iterated difference vanishes iff
$C(\zeta)=0$ for all $N$ roots $\zeta$ of $X^N-1$. If $\zeta$ has order $d<N$ then $d\mid r_i$ for some $i$
and $C(\zeta)=0$; if $\zeta$ has order $N$ no factor $\zeta^{r_i}-1$ vanishes and $C(\zeta)=0$ iff
$B(\zeta)=0$. So the condition is that $B$ vanishes at the primitive $N$-th roots of unity, which is
$\Phi_N\mid B(X)$.

For $N\mid M$ let $B_N(y)=\sum_{x\equiv y\ (N)}A(x)$; then $A(X)\equiv B_N(X)$ modulo $X^N-1$, so
$\Phi_N\mid A$ iff $\Phi_N\mid B_N$. For $N=p^a$ take $r_1=p^{a-1}$: $B_N(y)$ is the total of a class of depth
$a$, and $\nabla_{r_1}B_N=0$ says that the $p$ classes of depth $a$ inside each class of depth $a-1$ have equal
totals. For $N=q$ take $r_1=1$: $B_q$ is the vector of the root. For $N=p^aq$ take $r_1=N/p$, $r_2=N/q=p^a$;
identify $\Z_N$ with $\Z_{p^a}\times\Z_q$; subtracting $r_1$ moves a class of depth $a$ to the next one
inside the same class of depth $a-1$ (by an element of order $p$) and fixes the layer, and subtracting $r_2$
fixes the class and moves the layer by a generator of $\Z_q$. So $\nabla_{r_2}\nabla_{r_1}B_N=0$ says that
for neighbouring classes $c$, $c'$ of depth $a$ the vector $g(c)-g(c')$ is invariant under a generator of
$\Z_q$, that is, constant.
\end{proof}

So $E(A)=E(\mathrm{root})$, $F(A)=F(\mathrm{root})$, and $A$ is balanced iff the root is.

\subsection*{Pairs contained in a pattern} Let $N$ have $j$ levels below it. A pair $(E,F)$ of subsets of
$[1,j]$ is \emph{contained in the pattern of $N$} if $E\subseteq E(N)$ and $F\subseteq F(N)$. Its
\emph{shifted pair} is $E^-=\{l-1:l\in E,\ l\ge 2\}$, $F^-=\{l-1:l\in F,\ l\ge2\}$. Because the conditions
are ``for every class'',
\begin{equation}\label{eq:shift}
\parbox{0.84\textwidth}{if $(E,F)$ is contained in the pattern of $N$, then $(E^-,F^-)$ is contained in
the pattern of every child of $N$,}
\end{equation}
and $|E^-\cup F^-|$ is $|E\cup F|-1$ if $1\in E\cup F$ and $|E\cup F|$ otherwise. All lemmas below are stated
for pairs contained in a pattern, not for exact patterns; that is what makes the inductions work.

\section{Lower bounds}\label{sec:lower}

\subsection{Three sets of integers} For an integer $u\ge0$ put
\begin{gather*} S_u=\{p^ut+qm:\ t\ge1,\ m\ge0\},\qquad \Sp_u=\{p^ut+qm:\ t\ge1,\ m\ge1\},\\ Q_u=\{q\,p^ut:\ t\ge1\}. \end{gather*}
Each of the following is checked by writing out the definition; in (A4) and (A5) one uses that $q$ is a prime
other than $p$.
\begin{itemize}
\item[(A1)] For $u\ge1$: $p\,S_{u-1}\subseteq S_u$, $\ p\,\Sp_{u-1}\subseteq\Sp_u$, $\ p\,Q_{u-1}\subseteq Q_u$.
\item[(A2)] $S_u+S_u\subseteq S_u$; $\ S_u+q\Z_{\ge0}\subseteq S_u$; $\ \Sp_u+S_u\subseteq\Sp_u$;
 $\ \Sp_u+q\Z_{\ge0}\subseteq\Sp_u$; $\ S_u+q\Z_{\ge1}\subseteq\Sp_u$.
\item[(A3)] $Q_u+S_u\subseteq\Sp_u$ and $Q_u+q\Z_{\ge1}\subseteq\Sp_u$.
\item[(A4)] If $x\in S_u$ and $q\mid x$, then $x\in Q_u\cup\Sp_u$ and $x\ge q\,p^u$.
\item[(A5)] If $x\in\Sp_u$ and $q\mid x$, then $x\ge q\,p^u+q$.
\end{itemize}
(For (A4): $x=p^ut+qm$ with $q\mid p^ut$, so $q\mid t$ and $t\ge q$; if $m=0$ then $x\in Q_u$.)

\subsection{Two elementary lemmas}

\begin{lemma}[children of equal shapes]\label{lem:shapes}
Let $N$ be a class whose $p$ children have equal shapes, and $v$ the vector of a child of least total.
\begin{enumerate}
\item[(a)] The children have vectors $v+d_i\one$ with integers $d_i\ge0$, $|N|=p|v|+q\sum d_i$, and the
children have equal totals iff $\sum d_i=0$.
\item[(b)] $N$ is balanced iff every child is balanced. So if $N$ is unbalanced every child is unbalanced, in
particular not empty.
\item[(c)] If $A$ is balanced and the levels $1,\dots,a_0-1$ are all in $F(A)$, every class of depth
$\le a_0-1$ is balanced.
\end{enumerate}
\end{lemma}
\begin{proof}
(a) is the definition. (b) $g(N)=pv+(\sum d_i)\one$ is constant iff $v$ is, iff each $v+d_i\one$ is.
(c) By induction on the depth from the root, using (b).
\end{proof}

\begin{lemma}[divisibility]\label{lem:div}
If $E\subseteq E(N)$ then $p^{|E|}$ divides $|N|$; if moreover $N$ is balanced, $q\,p^{|E|}$ divides $|N|$.
\end{lemma}
\begin{proof}
Induction on the number of levels below $N$, by \eqref{eq:shift}: if $1\notin E$, $|N|$ is a sum of multiples
of $p^{|E|}$; if $1\in E$, the children have one total $s$, $p^{|E|-1}\mid s$ and $|N|=ps$. If $N$ is
balanced, $q\mid |N|$ as well.
\end{proof}

\subsection{The sizes of unbalanced classes}

\begin{lemma}\label{lem:S}
Let $N$ be an unbalanced class, $(E,F)$ a pair contained in its pattern, $u=|E\cup F|$. Then $|N|\in S_u$. In
particular $|N|\ge p^{|E\cup F|}$.
\end{lemma}
\begin{proof}
Induction on the number $j$ of levels below $N$; for $j=0$, $u=0$ and $S_0=\Z_{\ge1}$. Let $j\ge1$.

\emph{$1\in F$.} The children have equal shapes; by Lemma~\ref{lem:shapes} the child $v$ of least total is
unbalanced and $|N|=p|v|+q\sum d_i$. By induction $|v|\in S_{u-1}$, and $|N|\in S_u$ by (A1), (A2).

\emph{$1\in E\setminus F$.} The children have one total $s$, some child is unbalanced, $s\in S_{u-1}$ by
induction, and $|N|=ps\in S_u$.

\emph{$1\notin E\cup F$.} Some child is unbalanced, with total in $S_u$ by induction; every other child is
unbalanced with total in $S_u$ or balanced with total in $q\Z_{\ge0}$. Use (A2).
\end{proof}

Let $N$ be a class of depth $d$ and $1\le g\le n-d$. Say that $N$ \emph{has a witness at level $g$} if some
class $\Gamma$ of depth $d+g-1$ contained in $N$ has children of unequal totals (so $g\notin E(N)$). Such a
$\Gamma$ is not empty; if $g\ge2$ it lies in one child $N_c$ of $N$, and $N_c$ has a witness at its own level
$g-1$.

\begin{lemma}\label{lem:Sp}
Let $N$ be an unbalanced class, $(E,F)$ a pair contained in its pattern, $u=|E\cup F|$, and suppose $N$ has
a witness at some level $g\in F$. Then $|N|\in\Sp_u\cup Q_u$.
\end{lemma}
\begin{proof}
Induction on $j$; for $j=0$ there is no $g$. Let $j\ge1$ and $\Gamma$ the witness class.

\emph{$1\in F$.} Every child is unbalanced and $|N|=p|v|+q\sum d_i$ with $|v|\in S_{u-1}$
(Lemmas~\ref{lem:shapes}, \ref{lem:S}). If $\sum d_i\ge1$, $|N|\in\Sp_u$ by (A1), (A2). If $\sum d_i=0$ the
children of $N$ have equal totals, so $\Gamma\ne N$, $g\ge2$, and $\Gamma$ lies in a child $N_c$, unbalanced,
with a witness at its level $g-1\in F^-$; by induction $|N_c|\in\Sp_{u-1}\cup Q_{u-1}$ and
$|N|=p|N_c|\in\Sp_u\cup Q_u$.

\emph{$1\in E\setminus F$.} Then $g\ne1$, the children have one total $s$, and $\Gamma$ lies in a nonempty
child $N_c$ with a witness at its level $g-1\in F^-$. If $N_c$ is unbalanced, $s\in\Sp_{u-1}\cup Q_{u-1}$ by
induction. If $N_c$ is balanced, $q\mid s$; some child is unbalanced, so $s\in S_{u-1}$, and
$s\in Q_{u-1}\cup\Sp_{u-1}$ by (A4). Either way $|N|=ps\in\Sp_u\cup Q_u$.

\emph{$1\notin E\cup F$.} Then $g\ge2$ and $\Gamma$ lies in a nonempty child $N_c$ with a witness at its
level $g-1\in F^-$; here $|E^-\cup F^-|=u$. Let $R$ be the sum of the totals of the other children; each of
them is unbalanced with total in $S_u$ or balanced with total in $q\Z_{\ge0}$. If $N_c$ is unbalanced and
$|N_c|\in\Sp_u$: $|N|\in\Sp_u$ by (A2). If $N_c$ is unbalanced and $|N_c|\in Q_u$: if $R=0$ then
$|N|\in Q_u$; if some other child is unbalanced then $R\in S_u$ and $|N|\in\Sp_u$ by (A3); otherwise
$R\in q\Z_{\ge1}$ and $|N|\in\Sp_u$ by (A3). If $N_c$ is balanced, $|N_c|\in q\Z_{\ge1}$, some other child is
unbalanced, and $|N|\in\Sp_u$ by (A2).
\end{proof}

\begin{lemma}\label{lem:Spp}
In Lemma~\ref{lem:Sp} suppose moreover that no level $l<g$ lies in $E\setminus F$. Then $|N|\in\Sp_u$.
\end{lemma}
\begin{proof}
Induction on $j$. If $g\ge2$, the child $N_c$ containing $\Gamma$ has a witness at its level $g-1\in F^-$
and no level below $g-1$ lies in $E^-\setminus F^-$. \emph{$1\in F$}: as before
$|N|=p|v|+q\sum d_i$; if $\sum d_i\ge1$, $|N|\in\Sp_u$; if $\sum d_i=0$ then $g\ge2$, $N_c$ is unbalanced,
$|N_c|\in\Sp_{u-1}$ by induction and $|N|=p|N_c|\in\Sp_u$. \emph{$1\in E\setminus F$} does not occur
($g\in F$ gives $g\ge2$, and level $1<g$ would lie in $E\setminus F$). \emph{$1\notin E\cup F$}: $g\ge 2$; if
$N_c$ is unbalanced, $|N_c|\in\Sp_u$ by induction and $|N|\in\Sp_u$ by (A2); if $N_c$ is balanced,
$|N_c|\in q\Z_{\ge1}$, some other child is unbalanced with total in $S_u$, and $|N|\in\Sp_u$ by (A2).
\end{proof}

The hypothesis of Lemma~\ref{lem:Spp} cannot be dropped: for $p=2$, $q=3$ the class with children
$(6,0,0)=(3,0,0)+(3,0,0)$ and $(2,2,2)=(2,2,2)+(0,0,0)$ has $E=\{1\}$, $F=\{2\}$, $u=2$ and $12=q\,p^u$
elements, and $12\notin\Sp_2$. Its balanced child $(2,2,2)$ is the picture behind clause (4).

\subsection{The unbalanced half}

\begin{proposition}\label{prop:U}
Let $A$ be unbalanced, $E=E(A)$, $F=F(A)$, $e=|E|$, $k=|F\setminus E|$. Then $|A|\ge p^{e}$ if $k=0$, and
$|A|\ge p^{e}(p^{k}+q)$ if $k\ge1$.
\end{proposition}
\begin{proof}
Put $u=|E\cup F|=e+k$. By Lemma~\ref{lem:S}, $|A|\in S_u$, so $|A|\ge p^{u}$; this is the claim for $k=0$.
Let $k\ge1$ and $g\in F\setminus E$. Since $g\notin E$ the root has a witness at level $g$, and $g\in F$, so
$|A|\in\Sp_u\cup Q_u$ by Lemma~\ref{lem:Sp}. If $|A|=p^ut+qm$ with $t,m\ge1$: $p^e$ divides $|A|$
(Lemma~\ref{lem:div}) and $p^ut$, hence $m$, so $m\ge p^e$ and $|A|\ge p^u+q\,p^e$. If $|A|\in Q_u$:
$|A|\ge q\,p^{e+k}\ge p^e(p^k+q)$, because $(q-1)(p^k-1)\ge1$.
\end{proof}

\subsection{The balanced half}

\begin{lemma}\label{lem:bal}
Let $N$ be a balanced class whose $p$ children do not all have equal shapes, $(E,F)$ a pair contained in its
pattern, $u=|E\cup F|$. Then $|N|\ge q\,p^{u}$.
\end{lemma}
\begin{proof}
Some child is unbalanced, $q\mid|N|$, and $1\notin F$. If $1\in E$: the children have one total $s$, $q\mid s$,
$|E^-\cup F^-|=u-1$, $s\in S_{u-1}$ by Lemma~\ref{lem:S}, so $s\ge q\,p^{u-1}$ by (A4) and $|N|=ps$. If
$1\notin E$: $|E^-\cup F^-|=u$, the unbalanced children have totals in $S_u$ and the balanced ones in
$q\Z_{\ge0}$, so $|N|\in S_u$ by (A2) and $|N|\ge q\,p^u$ by (A4).
\end{proof}

\begin{proposition}\label{prop:B}
Let $A$ be a nonempty balanced multiset, $E=E(A)$, $F=F(A)$, $e=|E|$.
If $F=[1,n]$ then $|A|\ge q\,p^{e}$. If $F\ne[1,n]$, with $a_0$, $G$, $k'$, $x$ as in
Theorem~\ref{thm:main}(c), then $|A|\ge q\,p^{e}(p^{k'}+x)$.
\end{proposition}
\begin{proof}
The first statement is Lemma~\ref{lem:div} for the root. Let $F\ne[1,n]$.

\emph{The bound $q\,p^{e+k'}$.} Level $a_0\notin F$: some class $P$ of depth $a_0-1$ has children of unequal
shapes, and $P$ is balanced by Lemma~\ref{lem:shapes}(c). The pair
$E_P=\{a-a_0+1:a\in E,\ a\ge a_0\}$, $F_P=\{a-a_0+1:a\in F,\ a>a_0\}$ is contained in the pattern of $P$, and
$|E_P\cup F_P|=|E\cap[a_0,n]|+k'$; so $|P|\ge q\,p^{|E\cap[a_0,n]|+k'}$ by Lemma~\ref{lem:bal}. Going up from
$P$ to the root through the classes $P=P_{a_0-1}\subset\dots\subset P_0$, each step keeps or multiplies the
total, and multiplies it by $p$ at every level $d\in E$, $d<a_0$. Hence $|A|\ge q\,p^{e+k'}$.

\emph{The case $x=1$.} Then $k'\ge1$; put $g_1=\min G$. No level below $a_0$ is in $E$; and either
($\alpha$) $a_0\notin E$, or ($\beta$) $a_0\in E$, $p=2$, and no level of $E\setminus F$ lies strictly between
$a_0$ and $g_1$. The pair $(E_P,F_P)$ is contained in the pattern of every class of depth $a_0-1$; now
$|E_P|=e$ and $|E_P\cup F_P|=e+k'$.

\emph{Step 1.} If a class $P'\ne P$ of depth $a_0-1$ is not empty, it is balanced
(Lemma~\ref{lem:shapes}(c)), so $|P'|\ge q\,p^e$ by Lemma~\ref{lem:div}, and
$|A|\ge|P|+|P'|\ge q\,p^{e+k'}+q\,p^e$.

\emph{Step 2.} Otherwise $|A|=|P|$. As $g_1\in F\setminus E$, some class $\Gamma$ of depth $g_1-1\ge a_0$ has
children of unequal totals; it is not empty, so it lies in $P$, inside one child $N_c$ of $P$. By
\eqref{eq:shift} the pair $E^-=\{a-a_0:a\in E,\ a>a_0\}$, $F^-=\{a-a_0:a\in F,\ a>a_0\}$ is contained in the
pattern of each child of $P$; put $u^-=|E^-\cup F^-|=|E\cap(a_0,n]|+k'$. $N_c$ has a witness at its own level
$g=g_1-a_0\in F^-$.

\emph{Case ($\alpha$).} Here $|E^-|=e$ and $u^-=e+k'$, and some child of $P$ is unbalanced. If some child
$N_b$ of $P$ is balanced and not empty, $|N_b|\ge q\,p^e$ (Lemma~\ref{lem:div}); the sum $\Sigma$ of the totals
of the unbalanced children is in $S_{u^-}$ (Lemma~\ref{lem:S}, (A2)) and divisible by $q$, so
$\Sigma\ge q\,p^{u^-}$ by (A4), and $|A|\ge\Sigma+|N_b|$. Otherwise every nonempty child of $P$ is unbalanced;
then $N_c$ is, $|N_c|\in\Sp_{u^-}\cup Q_{u^-}$ by Lemma~\ref{lem:Sp}, and since $g(P)$ is constant and
$g(N_c)$ is not, some other child is unbalanced; let $R\ge1$ be the sum of the totals of the children other
than $N_c$, so $R\in S_{u^-}$. If $|N_c|\in Q_{u^-}$ then $|P|\ge q\,p^{u^-}+1$; if $|N_c|\in\Sp_{u^-}$ then
$|P|\in\Sp_{u^-}$ and $|P|\ge q\,p^{u^-}+q$ by (A5). Either way $|P|>q\,p^{e+k'}$; and $q\,p^e$ divides both
$|P|$ (Lemma~\ref{lem:div}) and $q\,p^{e+k'}$; so $|A|=|P|\ge q\,p^{e+k'}+q\,p^e$.

\emph{Case ($\beta$).} Here $|E^-|=e-1$ and $u^-=e-1+k'$. $P$ has two children $N_0,N_1$ of one total $s$,
$|P|=2s$, and $q$ is odd, so $q\mid s$. Their vectors sum to a constant, so if one were balanced both would
be, with equal shapes; hence both are unbalanced. $N_c$ is one of them, it has a witness at its level
$g\in F^-$, and no level $l<g$ lies in $E^-\setminus F^-$. By Lemma~\ref{lem:Spp}, $s\in\Sp_{u^-}$, so
$s\ge q\,p^{u^-}+q$ by (A5). As $q\,p^{e-1}$ divides $s$ and $q\,p^{u^-}$,
$s\ge q\,p^{u^-}+q\,p^{e-1}$ and $|A|=2s\ge q\,p^{e+k'}+q\,p^e$.
\end{proof}

The clauses enter as follows. (2): Step~1 needs that no level of $E$ lies above $a_0$; with one, a second
nonempty class of depth $a_0-1$ is forced only up to the factor $p$ it already contributes. (3): in case
($\beta$) with $p\ge3$ the child $N_c$ may be balanced, a third child being free to carry the witness. (4):
Lemma~\ref{lem:Spp} needs its hypothesis.

\section{Constructions}\label{sec:constructions}

Elements of $\Z_{p^n}$ are written as digit strings. For a set $T$ of levels let $D(T)$ be the set of
strings whose digits are $0$ outside $T$; for $b\notin T$ and a digit $c$ let $D(T;b\mapsto c)$ be the set of
strings with $u_b=c$ and $0$ at the other levels outside $T$. For a set $D$ of strings and $m\ge1$,
$m\cdot\mathrm{Fib}(D)$ is the multiset with $m$ copies of $(u,w)$ for every $u\in D$ and every layer $w$;
every class has a constant vector in it. $e_t$ is the unit vector of layer $t$.

\subsection{The unbalanced half} Given $E,F\subseteq[1,n]$, let $\tau(u)=1$ if $E\setminus F\ne\emptyset$
and every digit of $u$ at a level of $E\setminus F$ is nonzero, and $\tau(u)=0$ otherwise. Put
\[ A_1=\{(u,\tau(u)):u\in D(E\cup F)\},\qquad A_2=\mathrm{Fib}(D(E)), \]
and $A=A_1$ if $F\subseteq E$, $A=A_1+A_2$ otherwise. Then $|A|$ is $p^{e}$, respectively $p^{e+k}+q\,p^{e}$.

\begin{proposition}\label{prop:CU}
$A$ is unbalanced, $E(A)=E$ and $F(A)=F$.
\end{proposition}
\begin{proof}
Every class that meets $A$ meets $A_1$; its $A_1$-part has a power of $p$ elements, which $q$ does not
divide, and its $A_2$-part is a constant vector; so every nonempty class is unbalanced. At a level of $F$ the
$A_1$-parts of the children of a class have the same layer vector (the digit at that level does not enter
$\tau$) and the $A_2$-parts are constant: equal shapes. At a level of $E$ both parts split evenly: equal
totals. At a level of $F\setminus E$, in the class with all earlier digits $0$, the child with digit $0$
holds the whole $A_2$-part and the others none, while the $A_1$-parts are equal: unequal totals. At a level
outside $E\cup F$ a nonempty class has one nonempty child, and it is unbalanced: unequal totals and unequal
shapes. At a level $a\in E\setminus F$ take a class of depth $a-1$ that meets $A$ and whose digits at the
earlier levels of $E\setminus F$ are nonzero: the $A_1$-part of its child with digit $0$ lies in layer $0$,
that of the child with digit $1$ has the same size and reaches layer $1$, and the $A_2$-parts are constant;
so the two do not have equal shapes.
\end{proof}

\subsection{The balanced half} If $F=[1,n]$, $A=\mathrm{Fib}(D(E))$ has $q\,p^e$ elements, every class
balanced, equal totals exactly at the levels of $E$.

Let $F\ne[1,n]$, with $a_0$, $G$, $k'$ as in the theorem. The building block is a \emph{shape part}
$S=S(a_0,T_1,T_2,H,\text{mode})$, for $T_1\subseteq[1,a_0-1]$, $H\subseteq T_2\subseteq[a_0+1,n]$ and a mode
$\mathrm U$, $\mathrm Q$ or $\mathrm W$. Let $D^*$ be the strings that are $0$ outside
$T_1\cup\{a_0\}\cup T_2$ (the digit of level $a_0$ is free); $\tau(u)=1$ if $H\ne\emptyset$ and every digit
of $u$ at a level of $H$ is nonzero, and $0$ otherwise; and, with $(\Delta x)(w)=x(w)-x(w-1)$,
\begin{align*}
\text{mode U:}\quad&\Lambda_0x=x,\qquad \Lambda_1x=|x|\one-x;\\
\text{modes Q, W:}\quad&\Lambda_0x=|x|\one+\Delta x,\qquad\Lambda_1x=|x|\one-\Delta x;\\
\text{mode Q also:}\quad&\Lambda_ix=|x|\one\quad(2\le i\le p-1).
\end{align*}
$S$ has, for every $u\in D^*$ whose digit $i=u_{a_0}$ has a $\Lambda_i$ defined, and every layer $w$,
$(\Lambda_ie_{\tau(u)})(w)$ copies of $(u,w)$. (All these numbers are $\ge0$: the vectors are $e_t$,
$\one-e_t$, $\one+e_t-e_{t+1}$, $\one-e_t+e_{t+1}$, $\one$.) By linearity the vector of a class below
level $a_0$ with digit $i$ at level $a_0$ is $\Lambda_i\beta$, where $\beta$ counts the strings of $D^*$ in
the class with $\tau=0$ and with $\tau=1$ in layers $0$ and $1$; $\beta$ is supported on layers $0$ and $1$ and $|\beta|$ is a power of $p$ (odd, if $q=2$), so $\beta$ is
not constant, and $\Lambda_0,\Lambda_1$ map non-constant vectors to non-constant vectors. Summing over the
digit at level $a_0$ gives $|x|\one$, $p|x|\one$, $2|x|\one$ in modes U, Q, W. From this one reads off:

\begin{lemma}\label{lem:shape}
Let $\mu=1,p,2$ in modes $\mathrm U,\mathrm Q,\mathrm W$. Then $|S|=\mu\,q\,p^{|T_1|+|T_2|}$; every class of
depth $<a_0$ is balanced in $S$; and for every sum $\Phi$ of multisets $m\cdot\mathrm{Fib}(D)$ the multiset
$A=S+\Phi$ is balanced with $F(A)=[1,a_0-1]\cup(T_2\setminus H)$. In $S$: at a level of $T_1$ and at a level
of $T_2\setminus H$ the children of a class have equal vectors; at a level of $H$ they have equal totals;
at a level $a\ne a_0$ outside $T_1\cup T_2$ only the child with digit $0$ can be nonempty; at level $a_0$ the
children with digits $0,1$ of a class that meets $D^*$ have totals $p^{|T_2|}$ and $(q-1)p^{|T_2|}$ in mode
$\mathrm U$ and $q\,p^{|T_2|}$ each in modes $\mathrm Q$, $\mathrm W$, the remaining children having total
$q\,p^{|T_2|}$ in mode $\mathrm Q$ and $0$ in modes $\mathrm U$, $\mathrm W$.
\end{lemma}

Put $T_2^0=(E\cup F)\cap(a_0,n]$, $H^0=(E\setminus F)\cap(a_0,n]$, and mode$^0=\mathrm Q$ if $a_0\in E$,
$\mathrm U$ if not; so $T_2^0\setminus H^0=F\cap(a_0,n]$. Write $E_<=E\cap[1,a_0-1]$. Define $A$ by the first
case that applies:
\begin{itemize}
\item[(I)] $k'=0$: $A=S(a_0,E_<,T_2^0,H^0,\text{mode}^0)$.
\item[(III)] $k'\ge1$ and $E_<\ne\emptyset$; let $b=\min E$:
 \[A=S(a_0,E_<\setminus\{b\},T_2^0,H^0,\text{mode}^0)+\sum_{i=1}^{p-1}p^{k'}\cdot\mathrm{Fib}(D(E\setminus\{b\};b\mapsto i)).\]
\item[(IV)] $k'\ge1$, $a_0\in E$, $p\ge3$:
 \[A=S(a_0,E_<,T_2^0,H^0,\mathrm W)+\sum_{i=2}^{p-1}p^{k'}\cdot\mathrm{Fib}(D(E\setminus\{a_0\};a_0\mapsto i)).\]
\item[(V)] $k'\ge1$, $a_0\in E$, and some level of $E\setminus F$ lies strictly between $a_0$ and $\min G$;
 let $b$ be the least such:
 \[A=S(a_0,E_<,T_2^0\setminus\{b\},H^0\setminus\{b\},\mathrm Q)+\sum_{i=1}^{p-1}p^{k'}\cdot\mathrm{Fib}(D(E\setminus\{b\};b\mapsto i)).\]
\item[(II)] otherwise: $A=S(a_0,E_<,T_2^0,H^0,\text{mode}^0)+\mathrm{Fib}(D(E))$.
\end{itemize}
Cases (I), (III), (IV), (V) are clauses (1) to (4), and (II) is reached exactly when $x=1$.

\begin{proposition}\label{prop:CB}
In each case $A$ is balanced, $E(A)=E$, $F(A)=F$, and $|A|=q\,p^{e}(p^{k'}+x)$.
\end{proposition}
\begin{proof}[Outline of proof]
The verification level by level is written out in the working note (Section~6 of
\texttt{balanced-half/proof.md}) and is checked in Lean; here is its plan. $F(A)=F$ and balance are Lemma~\ref{lem:shape}, since in every case $T_2\setminus H=F\cap(a_0,n]$. Sizes: in
(I) and (II) the shape part has $q\,p^{e}$, resp.\ $q\,p^{e+k'}$, elements ($\mu=p$ accounts for $a_0\in E$),
and (II) adds $q\,p^e$. In (III), (IV), (V) one level $h$ of $E$ (the \emph{host level}: $b$, $a_0$, $b$) is
treated as if it were not in $E$ when the shape part is built, which then has $q\,p^{e-1+k'}$ elements in
(III), (V) and twice that in (IV); the fibre multisets fill the remaining $p-1$ (in (IV): $p-2$) children at
the host level with $q\,p^{e-1+k'}$ elements each, so the totals at the host level are equal and
$|A|=q\,p^{e+k'}$. A class at or below the host level contains either only elements of $S$ or only fibres of
one $\Phi_i$. For $E(A)$: at a level of $E$ other than $h$ both kinds of class split evenly
(Lemma~\ref{lem:shape}; fibres over $D(E\setminus\{h\};\dots)$ are free at the levels of $E\setminus\{h\}$).
At a level outside $E\cup F$, and at a level of $E^c$ before $a_0$, the class with all earlier digits $0$ is
nonempty and has one nonempty child. At $a_0\notin E$ (mode U) the children with digits $0$, $1$ of the class with all earlier digits $0$ have totals
$p^{|T_2|}$ and $(q-1)p^{|T_2|}$ in cases (I), (III), and $p^{c}(p^{k'}+q)$ and $p^{c}(q-1)p^{k'}$ with
$c=|E\cap(a_0,n]|$ in case (II), where the fibres lie in the first; any further child is empty. If $p\ge3$ a
nonempty child stands beside an empty one. If $p=2$ then $q$ is odd: in cases (I), (III) $q-1\ne1$, and in case
(II) one total is an odd multiple of $p^c$ and the other an even one.
At a level $g\in G$: in case (II) the class with all earlier digits $0$ has children with equal $S$-vectors
and all its fibres in the child with digit $0$; in cases (III), (IV), (V) a fibre-only class (digit $1$, $2$,
$1$ at the host level and $0$ elsewhere) is nonempty with all its elements in one child. In case (V) the
level $b$ itself is not in $F(A)$ by Lemma~\ref{lem:shape}, being outside $T_2$.
\end{proof}

Theorem~\ref{thm:main} is Propositions~\ref{prop:U}, \ref{prop:B}, \ref{prop:CU}, \ref{prop:CB} with
Lemma~\ref{lem:translate}.

\begin{remark}[one condition instead of four]
For $k'\ge1$, $x=0$ iff the pattern has a \emph{host level}: a level $b\in E$ with $b<\min G$ such that
$b<a_0$, or $b=a_0$ and $p\ge3$, or $b>a_0$, $b\notin F$ and $a_0\in E$. Equivalently (and without a list of
positions): $x=0$ iff some least multiset for the pattern $(E\cup G,F)$, which has $q\,p^{e+k'}$ elements,
contains a nonempty class of depth $<\min G$ all of whose subclasses are balanced. Concentrating the fibres
of such a class then makes the totals unequal at the levels of $G$ at no cost. The size condition cannot be
dropped. (Propositions~11 and~13 of the working note's proof.)
\end{remark}

\section{The formal proof}\label{sec:lean}

Theorem~\ref{thm:main} is formalised in Lean~4 (toolchain \lean{v4.37.0-rc1}) with Mathlib, in the library
\lean{Pnq} of the development that also holds the tiling theorems of the companion papers. %%WHERE%%

\subsection*{The statements} A multiset on $\Z_{p^nq}$ is a function \lean{w : ZMod (p \^{} n * q) → ℕ}, its
multiplicities. Its mask polynomial, with exponents the representatives in $[0,M)$, is
{\footnotesize
\begin{verbatim}
noncomputable def maskW (w : ZMod M → ℕ) : Polynomial ℤ :=
  ∑ a : ZMod M, Polynomial.C (w a : ℤ) * Polynomial.X ^ a.val
\end{verbatim}
}
and \lean{levelsEW w}, \lean{levelsFW w} are the sets of $\ell\in[1,n]$ with
\lean{Polynomial.cyclotomic (p \^{} ℓ) ℤ ∣ maskW w}, respectively
\lean{Polynomial.cyclotomic (p \^{} ℓ * q) ℤ ∣ maskW w}. The formulas of the theorem are four definitions:
{\footnotesize
\begin{verbatim}
def gOf (Es Fs : Finset ℕ) (a0 : ℕ) : Finset ℕ := (Fs \ Es).filter (fun ℓ => a0 < ℓ)

def xOf (p : ℕ) (Es Fs : Finset ℕ) (a0 : ℕ) : ℕ :=
  if gOf Es Fs a0 = ∅ ∨ (∃ ℓ ∈ Es, ℓ < a0) ∨ (a0 ∈ Es ∧ 3 ≤ p) ∨
      (a0 ∈ Es ∧ p = 2 ∧ ∃ b ∈ Es \ Fs, a0 < b ∧ ∀ g ∈ gOf Es Fs a0, b < g)
  then 0 else 1

def minU (p q : ℕ) (Es Fs : Finset ℕ) : ℕ :=
  if Fs ⊆ Es then p ^ Es.card else p ^ Es.card * (p ^ (Fs \ Es).card + q)

def minB (p q : ℕ) (Es Fs : Finset ℕ) (a0 : ℕ) : ℕ :=
  q * p ^ Es.card * (p ^ (gOf Es Fs a0).card + xOf p Es Fs a0)
\end{verbatim}
}
and its three parts are three theorems (\lean{Pnq/MinGroup.lean}):
{\footnotesize
\begin{verbatim}
theorem minsize_unbalanced_zmod (p q n : ℕ) [Fact p.Prime] [Fact q.Prime] (hpq : p ≠ q)
    (Es Fs : Finset ℕ) (hE : Es ⊆ Finset.Icc 1 n) (hF : Fs ⊆ Finset.Icc 1 n) :
    IsLeast {s | ∃ w : ZMod (p ^ n * q) → ℕ, ∑ a, w a = s ∧
        ¬ Polynomial.cyclotomic q ℤ ∣ maskW w ∧ levelsEW w = Es ∧ levelsFW w = Fs}
      (minU p q Es Fs)

theorem minsize_balanced_all_zmod (p q n : ℕ) [Fact p.Prime] [Fact q.Prime] (hpq : p ≠ q)
    (Es : Finset ℕ) (hE : Es ⊆ Finset.Icc 1 n) :
    IsLeast {s | ∃ w : ZMod (p ^ n * q) → ℕ, ∑ a, w a = s ∧ 0 < s ∧
        Polynomial.cyclotomic q ℤ ∣ maskW w ∧ levelsEW w = Es ∧
        levelsFW w = Finset.Icc 1 n}
      (q * p ^ Es.card)

theorem minsize_balanced_zmod (p q n : ℕ) [Fact p.Prime] [Fact q.Prime] (hpq : p ≠ q)
    (Es Fs : Finset ℕ) (hE : Es ⊆ Finset.Icc 1 n) (hF : Fs ⊆ Finset.Icc 1 n)
    (a0 : ℕ) (ha1 : 1 ≤ a0) (han : a0 ≤ n)
    (hlow : ∀ ℓ, 1 ≤ ℓ → ℓ < a0 → ℓ ∈ Fs) (hna : a0 ∉ Fs) :
    IsLeast {s | ∃ w : ZMod (p ^ n * q) → ℕ, ∑ a, w a = s ∧
        Polynomial.cyclotomic q ℤ ∣ maskW w ∧ levelsEW w = Es ∧ levelsFW w = Fs}
      (minB p q Es Fs a0)
\end{verbatim}
}
\noindent (line breaks differ from the files). \lean{IsLeast S m} is Mathlib's ``$m\in S$ and $m\le s$ for
every $s\in S$'', so each theorem has both halves of its part of Theorem~\ref{thm:main}: a multiset with exactly
the pattern and $m$ elements exists, and none with exactly the pattern has fewer. In the third, the
hypotheses on \lean{a0} say that it is the least level in $[1,n]$ outside \lean{Fs}. The three together cover
every pattern once.

\subsection*{The parts} \sloppy On the tree of Section~\ref{sec:tree} (\lean{Pnq/Tree.lean}: a $p$-ary tree of depth
$n$ whose leaves are functions \lean{Fin q → ℕ}, with $E$ and $F$ defined by recursion as in the text), the
lower bounds are \lean{pow\_card\_union\_le\_size}, \lean{sharp\_lemma3} (\lean{Pnq/Sharp.lean}) and
\lean{balanced\_base}, \lean{balanced\_plus} (\lean{Pnq/Balanced.lean}); the constructions are
\lean{exists\_unbalanced} (\lean{Pnq/BuildU.lean}) and \lean{exists\_balanced\_all}, \lean{exists\_balanced}
(\lean{Pnq/BuildB.lean} with five files of lemmas); \lean{Pnq/MinTree.lean} joins them into three
\lean{IsLeast} statements about trees. All of these hold for coprime $p,q\ge2$, prime or not.
\lean{Pnq/Weighted.lean} is the bridge to the group: the Fourier zeros of a multiset, their equivalence with
cyclotomic divisibility of the mask polynomial (\lean{zeroAtW\_iff\_cyclotomic\_dvd}), the tree
\lean{treeOfW w} of a multiset, the translation of Lemma~\ref{lem:translate} (\lean{size\_treeOfW},
\lean{E\_treeOfW}, \lean{F\_treeOfW\_zero}, \lean{F\_treeOfW}), and that every tree is the tree of a
multiset (\lean{treeOfW\_surjective}; the Chinese remainder theorem); it also proves that on a set these
notions are the ones used for the tiling theorems. The new files have $3{,}478$ lines.

\subsection*{How the formal proof differs from this paper} The formal lower bound for part (a) is not
Proposition~\ref{prop:U}: it is the direct induction of the working note (two claims and four cases), written
before the argument by sets of sizes was found. The formal lower bounds for part (c) follow
Lemmas~\ref{lem:shapes} to~\ref{lem:bal} and Proposition~\ref{prop:B}. The formal constructions are recursive
builders of trees that follow Section~\ref{sec:constructions}, with changes of convenience that the instances
that wrote them reported: a string counts as having all its digits at the levels of $H$ (or of
$E\setminus F$) nonzero also when there are no such levels, which renames layers; in modes $\mathrm Q$,
$\mathrm W$ the vector for digit $0$ is $\one+e_0-e_2$ for $q\ge3$ and $\one+e_1-e_0$ for $q=2$; and in case
(III) any level of $E$ before $a_0$ serves as host. The remark on host levels at the end of
Section~\ref{sec:constructions} is not formalised.

\subsection*{How it was made, and what was checked} The author fixed every definition and every theorem
statement first, in files that compiled with the proofs left open, and before any proof was written had Lean
evaluate \lean{minU}, \lean{minB} and $q\,p^{e}$ on the $1{,}265$ patterns whose least size the solvers had
decided (Section~\ref{sec:computations}): all agree. Four fresh instances of the model then each wrote the
proofs of one part in a copy of its own, forbidden to change a statement. On the merged copy the author
checked that every fixed declaration stands character for character; that Lean prints each fixed definition
and statement exactly as it printed them before the proofs existed; that the new files contain no
\lean{sorry}, no added axiom, and no instance, notation, macro or attribute that could change what a statement
means; that the whole development builds; and that \lean{\#print axioms} gives, for each of the theorems named
above, only \lean{propext}, \lean{Classical.choice} and \lean{Quot.sound}. %%NOON%% %%AUDIT%%

\subsection*{What the formal proof does not show} That the theorem is new; that the prose proofs of this
paper are right (they were read by the author and, in their working-note form, refereed by fresh instances,
which is a weaker thing); anything about sets; and nothing has been read by a person.


\section{Computations}\label{sec:computations}

The formulas were found from, and are confirmed by, exhaustive computation. For each pattern the least size
was computed by two integer-programming models made independently of one another (CP-SAT on class counts,
and HiGHS on polynomial remainders modulo $\Phi_d$), every witness being checked a third way by complex
values: all $2^{2n+1}$ patterns for $n=2$ and $(p,q)\in\{(2,3),(3,2),(2,5),(5,2),(3,5),(5,3)\}$, for $n=3$
and several of these pairs, and for $n=4$, $(p,q)=(2,3)$, as far as the solvers finished: $1{,}265$ patterns
decided in all, $1{,}181$ of them by both solvers. Theorem~\ref{thm:main} gives the computed value in every
one. The rule of (c) was fitted on $n\le3$ and sealed before the $n=4$ results were opened; an earlier guess
of the author's, without clauses (2) and (4), was wrong on $39$ of $554$ patterns. The floating-point solver
is a verdict and not a proof; the formal proof makes it unnecessary, and the table is now a check on the
formal \emph{statement}: Lean's own evaluation of the formula agrees with all $1{,}265$ computed values
(Section~\ref{sec:lean}).

\begin{thebibliography}{9}
\bibitem{CM} E.~Coven and A.~Meyerowitz, \emph{Tiling the integers with translates of one finite set},
 J.~Algebra 212 (1999), 161--174.
\bibitem{KLaMS} G.~Kiss, I.~{\L}aba, C.~Marshall and G.~Somlai, \emph{Lower bounds for mask polynomials with many
cyclotomic divisors}, Adv. Math. 494 (2026), 110932; arXiv:2507.11672. (Read in part by the author; read
 through by a fresh instance of the model.)
\bibitem{KMSV} G.~Kiss, R.~D.~Malikiosis, G.~Somlai and M.~Vizer, \emph{On the discrete Fuglede and Pompeiu
 problems}, Anal. PDE 13 (2020), 765--788; arXiv:1807.02844. (Proposition~3.8 and Theorem~4.9 were read by a
 fresh instance of the model; the author has not read the paper.)
\bibitem{LL} T.~Y.~Lam and K.~H.~Leung, \emph{On vanishing sums of roots of unity}, J.~Algebra 224 (2000),
 91--109; arXiv:math/9511209. (Read in part by a fresh instance of the model; not by the author.)
\bibitem{LaM} I.~{\L}aba and C.~Marshall, \emph{Vanishing sums of roots of unity and the Favard length of
 self-similar product sets}, Discrete Analysis 2022:19; arXiv:2202.07555. (Theorem~1.2, Corollary~6.4 and
 their proofs were read by fresh instances of the model; the author has not read the paper.)
\bibitem{Mal} R.~D.~Malikiosis, \emph{On the structure of spectral and tiling subsets of cyclic groups}, Forum
 Math. Sigma 10 (2022), e23; arXiv:2005.05800. (The proof of Theorem~6.1 was read by a fresh instance of the
 model and by Noon; the author checked the lines derived from it.)
\bibitem{MK} R.~D.~Malikiosis and M.~N.~Kolountzakis, \emph{Fuglede's conjecture on cyclic groups of order
 $p^nq$}, Discrete Analysis 2017:12; arXiv:1612.01328. (The author read pp.~10--11.)
\end{thebibliography}

\end{document}
